% 金星（综述）
% license CCBYSA3
% type Wiki

本文根据 CC-BY-SA 协议转载翻译自维基百科\href{https://en.wikipedia.org/wiki/Venus}{相关文章}。

金星是距离太阳第二近的行星。由于其轨道最接近地球，且两者同为类地岩石行星，并在大小、质量和表面引力上几乎完全一致，因此金星常被称为地球的“孪生”或“姐妹”行星。然而，金星在许多方面与地球显著不同，尤其是它没有液态水，其大气层远比太阳系中其他任何岩石天体都更加厚重和致密。金星大气主要由二氧化碳组成，并覆盖着一层浓密的硫酸云层，环绕整个行星。在平均地表高度处，大气温度达到 737\,K（464\,°C；867\,°F），压力为地球海平面的 92 倍，使其最低层大气呈现超临界流体状态。从地球观测时，金星呈现为类似恒星的亮点，比天空中任何其他天然光源都更明亮[22][23]，并且作为一颗内行星，它总是出现在接近太阳的方向，要么作为最明亮的“启明星”，要么作为“长庚星”。

金星与地球的轨道使两者在约 1.6 年的会合周期中彼此接近。在此过程中，金星会比其他任何行星更接近地球；相较之下，由于水星的轨道更靠近太阳，因此它在自身轨道路径中整体上比其他行星更接近地球。在从地球出发的行星际空间飞行中，金星常被用作引力助推的中转点，从而提供更快速且更经济的航程。金星没有卫星，并以非常缓慢的逆行方式绕其自转轴旋转，这被认为是太阳潮汐锁定效应与金星巨大大气层的受热差异相互竞争的结果。因此，金星的一昼夜长达 116.75 个地球日，而其一个太阳年为 224.7 个地球日。

金星具有微弱的磁层；由于缺乏内部发电机，其磁场是由太阳风与大气相互作用所诱导产生的。在内部结构上，金星具有核心、地幔和地壳。内部热量通过活跃的火山活动逸出[24][25]，从而导致地表更新，而不是板块构造。金星在其早期历史中可能曾拥有液态地表水及宜居环境[26][27]，但由于失控的温室效应，所有水分最终蒸发，使金星演变为当前的状态[28][29][30]。在云层高度处的大气条件是整个太阳系中与地球最为相似的，这些条件被认为可能有利于金星生命的存在，并在 2020 年发现了潜在的生物标志物，从而推动了新的研究与探测任务。

在人类历史上，全世界的人们都曾观测过金星，并使其在众多文化中具有特殊的重要性。随着望远镜的出现，金星的相位变得可辨识，并在 1613 年被呈现为推翻当时占主导地位的地心模型、支持日心模型的决定性证据。1961 年，“金星 1 号”首次造访金星，飞掠行星并实现了人类第一次行星际飞行。1962 年，第二次行星际任务“水手 2 号”带回了首批来自金星的数据。1967 年，首个行星际撞击器“金星 4 号”抵达金星，随后在 1970 年由着陆器“金星 7 号”成功登陆。截至 2025 年，“太阳轨道飞行器”正前往于 2026 年飞掠金星，而下一项计划发射前往金星的任务是“金星生命探测器”，预定同样于 2026 年发射。
\subsection{物理特征}
\begin{figure}[ht]
\centering
\includegraphics[width=8cm]{./figures/23fb842f59355847.png}
\caption{金星（从左数第二，伪彩色显示）按比例与太阳系内侧具有行星质量的天体一起展示，排列顺序依其轨道自太阳向外依次排列（从左到右：水星、金星、地球、月球、火星和谷神星）。} \label{fig_Venus_1}
\end{figure}
\begin{figure}[ht]
\centering
\includegraphics[width=8cm]{./figures/f163eb810133acdf.png}
\caption{金星在不同波长下的成像} \label{fig_Venus_2}
\end{figure}
金星是太阳系中四颗类地行星之一，这意味着它与地球一样属于岩石天体。它在大小和质量上与地球相似，因此常被描述为地球的“姐妹”或“孪生”行星[31]。由于其自转缓慢，金星的形状非常接近球形[32]。其直径为 12,103.6\,km（7,520.8\,mi），仅比地球小 638.4\,km（396.7\,mi）；其质量为地球的 81.5\%，使金星成为太阳系中第三小的行星。金星表面的环境与地球完全不同，因为其致密大气中 96.5\% 为二氧化碳，导致强烈的温室效应，其余约 3.5\% 为氮气[33]。金星表面气压为 9.3\,MPa（93\,bar），平均表面温度为 737\,K（464\,°C；867\,°F），高于两种主要组分的临界点，使其地表大气成为一种超临界流体，主要由超临界二氧化碳及部分超临界氮组成。
\subsubsection{自然历史}
\textbf{形成}

包括金星在内的类地岩石行星被认为经历了五个阶段形成：尘埃沉降、微行星形成、行星胚胎阶段、巨撞阶段，以及最终的大气形成阶段。来自金星的有限测量数据使其形成时间线难以进行更详细的分析[34]。

\textbf{未来}

金星预计将在七至八十亿年后当太阳演变为红巨星时，与水星一起被毁灭，也有可能包括地球和月球[35]。
\subsubsection{地理}
\begin{figure}[ht]
\centering
\includegraphics[width=8cm]{./figures/884bdec4a9f38b8d.png}
\caption{彩色编码的高程地图，显示金星上呈黄色的高地“大陆”以及次要地形特征} \label{fig_Venus_3}
\end{figure}
\begin{figure}[ht]
\centering
\includegraphics[width=8cm]{./figures/4ef3c5c0b6c896cf.png}
\caption{金星表面雷达数据的球形视图，突出显示地表特征（1989 年，麦哲伦号）。这些颜色并不代表地表的真实外观[36]。} \label{fig_Venus_4}
\end{figure}
在 20 世纪探测器揭示部分秘密之前，金星表面一直是推测的对象。1975 年和 1982 年的“金星号”着陆器传回了由沉积物和相对棱角分明的岩石覆盖的地表图像[37]。1990–91 年间，“麦哲伦号”对其表面进行了详细测绘。有大量火山活动的证据，并且大气中二氧化硫的变化可能表明存在活火山[38][39]。

大约 80\% 的金星表面被光滑的火山平原覆盖，其中 70\% 为带有皱褶脊的平原，10\% 为光滑或叶状平原[40]。其余的地表由两个高地“大陆”构成，一个位于行星的北半球，另一个位于赤道以南。北部高地称为伊什塔尔高原（Ishtar Terra），得名于巴比伦的爱神伊什塔尔，其面积约与澳大利亚相当。麦克斯韦山脉位于伊什塔尔高原，其最高峰斯卡迪山（Skadi Mons）是金星上的最高点，海拔比金星平均表面高度高 11\,km（7\,mi）[41]。南部高地称为阿芙洛狄忒高原（Aphrodite Terra），得名于希腊神话中的爱神，是两大高地中较大的一个，面积约与南美洲相当。该区域的大部分被裂缝和断层网络覆盖[42]。

近期有关于金星上岩浆流动的证据（2024 年）[43]，例如在盾状火山 Sif Mons 的熔岩流以及在平坦平原 Niobe Planitia 上的熔岩流[24]。地表可见破火山口。金星的撞击坑数量很少，表明其表面相对年轻，约为 3 亿至 6 亿年[44][45]。金星除了拥有岩石行星上常见的撞击坑、山脉和山谷外，还具有一些独特的地表特征。其中包括称为“farra”的平顶火山构造，其形状类似煎饼，直径从 20 至 50\,km（12 至 31\,mi）不等，高度在 100 至 1,000\,m（330 至 3,280\,ft）之间；称为“novae”的放射状、星形裂缝系统；具有放射状和同心裂隙、类似蜘蛛网的构造，被称为“arachnoids”；以及“coronae”，即被洼地包围的环状裂缝结构。这些特征均具有火山成因[46]。
\begin{figure}[ht]
\centering
\includegraphics[width=8cm]{./figures/5421d90f38272eab.png}
\caption{彩色化影像（“金星 9 号”，1975 年），金星天空在地表呈橙黄色，这是由于瑞利散射或低层大气中的蓝色吸收体所致，而在更高的高度则呈白色[47][48]；同时其地表为类似玄武岩的深灰色，可能因氧化而呈红色[36]。} \label{fig_Venus_5}
\end{figure}
大多数金星表面特征以历史和神话中的女性命名[49]。例外包括麦克斯韦山脉，以詹姆斯·克拉克·麦克斯韦命名，以及阿尔法高地区（Alpha Regio）、贝塔高地区（Beta Regio）和奥夫达高地区（Ovda Regio）。这后三个特征是在国际天文学联合会（负责行星命名的机构）采用现行系统之前命名的[50]。

金星地形特征的经度以其本初子午线为基准表达。最初的本初子午线通过位于阿尔法高地区南部、椭圆形地形“夏娃”（Eve）中央的雷达亮点[51]。在“金星号”任务完成后，本初子午线被重新定义为通过塞德娜平原（Sedna Planitia）上阿里阿德涅（Ariadne）陨石坑中央峰的位置[9][52]。

地层学上最古老的镶嵌地形（tessera terrains）在“金星快车号”和“麦哲伦号”测量中显示出相对于周围玄武质平原更低的热辐射率，这表明其具有不同的、可能更富长英质（felsic）的矿物组合[29][53]。生成大量长英质地壳的机制通常需要存在海洋和板块构造，这意味着早期金星上可能存在宜居条件，并在某一时期拥有大量水体[54]。然而，镶嵌地形的具体性质仍不确定[55]。

2023 年的研究首次提出金星在古代可能存在板块构造，因此可能拥有更宜居的环境，甚至可能具备维持生命的能力[26][27]。金星因此成为研究类地行星形成及其宜居性的重要案例。

\textbf{火山活动}
\begin{figure}[ht]
\centering
\includegraphics[width=8cm]{./figures/96e9067c358877bf.png}
\caption{金星埃斯特拉区域内两个煎饼穹丘的雷达拼接图——二者均宽 65\,km（40\,mi），高度不足 1\,km（0.62\,mi）} \label{fig_Venus_6}
\end{figure}


